\section{Proofs and asymptotic calculations}
\label{app:proofs}

\subsection{The volatility-space density factor}
\label{app:durrlemanproof}

Let $c(k)=\Black(k,w(k))$.  The identity
$e^k\phin(d_-)=\phin(d_+)$ gives
\[
\partial_k\Black=-e^k\PhiN(d_-),
\qquad
\partial_w\Black=\frac{\phin(d_+)}{2\sqrt w}.
\]
The required second partial derivatives are
\begin{align*}
\partial_{kk}\Black
&=-e^k\PhiN(d_-)+\frac{\phin(d_+)}{\sqrt w},\\
\partial_{kw}\Black
&=\frac{d_+\phin(d_+)}{2w},\\
\partial_{ww}\Black
&=-\frac{d_+\phin(d_+)}{2\sqrt w}\,\partial_wd_+
-\frac{\phin(d_+)}{4w^{3/2}},
\qquad
\partial_wd_+=\frac{k}{2w^{3/2}}+\frac1{4\sqrt w}.
\end{align*}
Twice differentiating $c(k)$ and subtracting the first derivative gives
\[
c''-c'=(\partial_{kk}\Black-\partial_k\Black)
+(2\partial_{kw}\Black-\partial_w\Black)w'
+\partial_{ww}\Black(w')^2+\partial_w\Black w''.
\]
Now
\[
2\partial_{kw}\Black-\partial_w\Black
=-\frac{k\phin(d_+)}{w^{3/2}},
\qquad
d_+\partial_wd_+=-\frac{k^2}{2w^2}+\frac18.
\]
Factoring $\phin(d_+)/\sqrt w$ yields
\[
c''-c'=\frac{\phin(d_+)}{\sqrt w}
\left[
\left(1-\frac{kw'}{2w}\right)^2
-\frac{(w')^2}{4}\left(\frac1w+\frac14\right)
+\frac{w''}{2}\right].
\]
Finally $f_X(k)=e^{-k}(c''-c')$ and
$e^{-k}\phin(d_+)=\phin(d_-)$, which gives \eqref{eq:durrleman}.

\subsection{The ATM identities}
\label{app:atmproof}

At $(k,w)=(0,w_0)$, with $d=\sqrt{w_0}/2$,
\begin{equation}
\begin{aligned}
\partial_k\Black&=-(1-\PhiN(d)),
&\partial_w\Black&=\frac{\phin(d)}{2\sqrt{w_0}},\\
\partial_{kk}\Black&=-(1-\PhiN(d))
+\frac{\phin(d)}{\sqrt{w_0}},
&\partial_{kw}\Black&=\frac{\phin(d)}{4\sqrt{w_0}},\\
\partial_{ww}\Black&=-\frac{\phin(d)}{4w_0^{3/2}}
-\frac{\phin(d)}{16\sqrt{w_0}}.
\end{aligned}
\label{eq:blackpartials}
\end{equation}
From \eqref{eq:cprime}--\eqref{eq:bl_division},
\[
c'(0)=-(1-u_*),\qquad c''(0)=f_*-(1-u_*).
\]
One derivative of $\Black(k,w(k))=c(k)$ gives
\[
w'(0)=
\frac{-(1-u_*)+(1-\PhiN(d))}{\phin(d)/(2\sqrt{w_0})}
=\frac{2\sqrt{w_0}}{\phin(d)}\Delta.
\]
Two derivatives give
\[
\partial_{kk}\Black+2\partial_{kw}\Black w'
+\partial_{ww}\Black(w')^2+\partial_w\Black w''=c''.
\]
Substitution of \eqref{eq:blackpartials} and the first-derivative formula
reduces this to \eqref{eq:atm_wpp}.  Finally,
\[
\sigma'=\frac{w'}{2\sqrt{\tau w}},
\qquad
\sigma''=\frac{w''}{2\sqrt{\tau w}}
-\frac{(w')^2}{4\sqrt\tau\,w^{3/2}},
\]
which gives \eqref{eq:atm_sigma1}--\eqref{eq:atm_sigma}.

\subsection{Sublinear wing asymptotics}
\label{app:sublinearproof}

Consider the right wing with $0<\alpha_+\le1/2$ and set
\[
p=\frac1{1-\alpha_+},
\qquad
A=\left(\frac{1-\alpha_+}{\lambda_+}\right)^p.
\]
Proposition \ref{prop:tails} gives
$-\log\Prob(X>x)\sim Ax^p$.  The call has the tail-integral form
\[
c(k)=\int_k^\infty e^x\Prob(X>x)\,\dd x.
\]
Since $p>1$, endpoint Laplace bounds give
\begin{equation}
-\log c(k)\sim Ak^p.
\label{eq:calllogtail}
\end{equation}
For example, for every $\varepsilon>0$ and all large $x$,
$e^{-(A+\varepsilon)x^p}\le\Prob(X>x)
\le e^{-(A-\varepsilon)x^p}$; integrating the two bounds and then sending
$\varepsilon$ to zero proves \eqref{eq:calllogtail}.

All positive moments exist, so Lee's formula gives $w(k)/k\to0$.  Mills'
ratio in the Black formula then gives
\begin{equation}
-\log\Black(k,w(k))
\sim\frac{k^2}{2w(k)}.
\label{eq:blacklogtail}
\end{equation}
Equating \eqref{eq:calllogtail} and \eqref{eq:blacklogtail} yields
\[
w(k)\sim\frac1{2A}k^{2-p},
\]
which is \eqref{eq:rightsublinearwing}.  The left proof uses the put and is
identical.

\subsection{Approximation at fixed tail class}
\label{app:universalproof}

Let $g^*$ satisfy \cref{prop:universal}.  By Weierstrass' theorem there are
polynomials $g_N$ with
$\varepsilon_N=\lVert g_N-g^*\rVert_\infty\to0$.  The endpoint chord and
$P_2,\ldots,P_N$ span every polynomial of degree at most $N$.

The speed ratio is bounded uniformly on the whole line:
\[
e^{-\varepsilon_N}
\le\frac{x_N'(z)}{x^{*\prime}(z)}
\le e^{\varepsilon_N}.
\]
Consequently $\bar x_N\to\bar x^*$ uniformly on every compact $z$ interval.
It remains to control the tails and the martingale shifts.

If $\alpha_+>0$, boundedness of the $g_N$ gives constants $C_1,C_2$ such that
\[
\bar x_N(z)\le C_1+C_2z^{1-\alpha_+},\qquad z\ge0.
\]
Hence $e^{\bar x_N(z)}\rho(z)$ is dominated by an integrable multiple of
$\exp(C_2z^{1-\alpha_+}-z)$.  If $\alpha_+=0$, the strict endpoint margin
$e^{g^*(1)}<1$ supplies $\bar\lambda<1$ and a common bound
$e^{\bar x_N(z)}\rho(z)\le C e^{-(1-\bar\lambda)z}$ for large $z$.
On the left, $\bar x_N(z)\le0$ and $\rho(z)\le e^z$.  Dominated convergence
therefore gives convergence of the normalizers and shifts.

The quantiles converge at every fixed rank.  The same envelopes provide a
common moment of order strictly larger than one, so the gross returns are
uniformly integrable.  Under the common uniform-rank coupling,
\[
\sup_k|c_N(k)-c^*(k)|
\le\E|e^{Q_N(U)}-e^{Q^*(U)}|\longrightarrow0.
\]
This proves uniform convergence of call prices.

\subsection{Tail restrictions imposed by calendar order}
\label{app:calendartails}

On the right, \eqref{eq:righttransporttail} and endpoint integration give
\[
\log G(z)=-z+
\frac{\lambda_+}{1-\alpha_+}z^{1-\alpha_+}+o(z^{1-\alpha_+}).
\]
If a farther expiry has a larger $\alpha_+$, its positive correction has a
smaller power and its ledger is eventually below the nearer ledger.  If the
exponents are equal, the larger $\lambda_+$ gives the larger eventual ledger.
The left lower-share ledger satisfies
\[
\log(1-G(z))=z-
\frac{\lambda_-}{1-\alpha_-}|z|^{1-\alpha_-}
+o(|z|^{1-\alpha_-}),
\qquad z\to-\infty.
\]
The same comparison gives the left conditions stated after
\eqref{eq:tailscalecalendar}.  Equal exponents and equal endpoint coefficients
leave the additive asymptotic constants to decide the order.
